% ═══════════════════════════════════════════════════════════════════════════
\chapter{Problema a Trabajar}
% ═══════════════════════════════════════════════════════════════════════════

Sea $m_t$ el precio medio del activo y $X_t \in \mathbb{R}^d$ un vector de variables
de mercado observables en el instante $t$ (spread actual, desequilibrio del libro de
órdenes, volatilidad realizada y volumen en ventanas causales). Se denotan por
$\lambda_0^a(t\mid X_t)$ y $\lambda_0^b(t\mid X_t)$ las intensidades base de trades
agresores que pueden alcanzar, respectivamente, una cotización pasiva en el ask o en
el bid. La intensidad de fill a una distancia elegida $\delta$ combina esa intensidad
base con una función decreciente de la distancia y una regla conservadora de ejecución
reconstruida mediante replay del LOB.

El problema central de este trabajo se articula en cuatro preguntas de investigación:

\begin{enumerate}[label=\textbf{P\arabic*.}]
  \item \textquestiondown{}Cómo modelar $\lambda_0^a(t\mid X_t)$ y $\lambda_0^b(t\mid X_t)$ como
        funciones del estado del mercado $X_t$ y del tiempo $t$ usando un enfoque
        semiparamétrico que preserve la interpretabilidad del modelo de control
        estocástico?
  \item \textquestiondown{}Qué diferencias en PnL neto y control de inventario se obtienen frente a
        Avellaneda--Stoikov y frente a una política basada en intensidades Hawkes?
  \item \textquestiondown{}Cómo impacta la estimación data-driven de las intensidades en la política
        óptima de cotización derivada vía la ecuación HJB?
  \item \textquestiondown{}El desempeño del estimador de $g(X_t)$ se mantiene fuera de muestra ante
        cambios de régimen, costos, latencia y especificaciones de inventario?
\end{enumerate}

